\chapter{结论与展望}

\section{研究工作总结}

本文围绕某二拖一燃气-蒸汽联合循环热电联产机组，基于历史运行数据建立了燃气轮机、余热锅炉、汽轮机以及联合循环整机的性能分析模型，并从能量分析和相关性分析两个层面对机组运行特性进行了研究。主要工作和结论如下。

首先，本文建立了统一的工质流股建模方法。针对联合循环机组中空气、天然气、烟气、水和水蒸气等多类工质并存的问题，本文将工质状态统一表示为包含温度、压力、质量流量、组分、比焓、比熵和比㶲等参数的流股对象。该方法使不同设备之间的质量流和能量流能够连续传递，为后续部件模型和整机模型提供了统一的数据基础。

其次，本文建立了燃气轮机部件级性能分析模型。将压气机、燃烧室和透平分别作为控制体，计算压气机耗功和等熵效率，燃烧室燃料热输入、出口烟气状态，以及透平输出功、等熵效率和等熵损失功率。结果表明，各个计算值具有较高的精度，能够较好地反映燃气轮机的运行状态和性能水平。

再次，本文建立了余热锅炉和汽轮机的汽水侧性能分析模型。余热锅炉模型从整台设备的烟气侧放热和汽水侧吸热出发，计算换热效率、烟气能量释放、水侧吸热分配以及不同压力等级蒸汽的热量贡献。汽轮机模型将高压缸、中压缸和低压缸分别建模，计算各缸输出功率和等熵效率，并考虑供热工况下低压蒸汽进入热网导致的通流方式变化。结果表明，余热锅炉高压蒸汽吸热和汽轮机高、中压缸出力是底循环发电能力的重要来源；供热工况下，部分蒸汽能量由发电侧转移至热网侧，使汽轮机出力和热电效率之间呈现不同于非供热工况的变化规律。

最后，本文对联合循环效率影响因素进行了相关性分析。非供热工况下，联合循环发电效率与汽机出力、整机负荷、主蒸汽参数和余热回收能力密切相关，说明提高燃气轮机高效运行水平和增强蒸汽循环做功能力是提升发电效率的关键。供热工况下，热电效率与供热量、环境条件和燃机总出力关系更为密切，表明供热需求和季节性环境变化会显著改变机组能量分配方式。相关性分析还表明，多数运行参数之间存在明显耦合关系，效率优化需要结合部件热力模型进行综合判断，不能仅依据单一变量变化作出结论。

总体而言，本文建立的模型能够基于运行数据对联合循环机组进行部件级和整机级性能评价，较好地反映燃气轮机、余热锅炉、汽轮机和供热系统之间的能量耦合关系。研究结果可为热电联产机组的运行诊断、性能评价和运行优化提供参考，并有望应用到实际电厂的在线性能监测和运行优化系统中。

\section{未来工作展望}

本文研究仍存在一些不足，后续可从以下几个方面进一步完善。

第一，空气流量测点的偏差较大，对后续压气机和透平的性能计算也产生一定影响，需要寻找更可靠的修正方法。后续可进一步引入环境参数、负荷水平和设备运行状态等变量，建立更具泛化能力的空气流量修正模型；同时，可结合燃气轮机功率平衡关系反算空气流量，以提高缺失测点条件下质量流量边界的可靠性。

第二，可进一步细化余热锅炉和汽轮机模型。本文仅从整台余热锅炉的能量输入输出角度分析换热性能，尚未对省煤器、蒸发器、过热器和再热器等各段换热器分别建模，也未充分考虑给水泵、汽包、减温器和节流阀等辅助设备对系统性能的影响。初步分析表明，减温器混合过程引起的不可逆损失较为显著，通过优化减温水调节策略有望提高整机运行效率。后续可建立更精细的分段换热模型，分析不同工况下各换热器的传热特性及其对机组性能的影响，为设备改造和运行优化提供更具针对性的指导。

第三，影响因素分析有待深化。本文采用 Pearson 相关系数识别了与效率变化关联最强的运行参数，但尚未深入区分各变量之间的耦合关系。后续仍需结合部件级热力模型和多变量分析方法，进一步辨识关键参数对效率变化的独立影响及其可调节路径。

第四，可开展面向运行优化的工况寻优研究。本文主要完成性能计算和影响因素分析，后续可在此基础上建立以发电效率、热电效率、供热量或燃料消耗量为目标的优化模型，研究燃气轮机负荷分配、供热抽汽量、主蒸汽参数和余热锅炉运行边界的协同优化方法，为实际运行调度提供更直接的决策支持。

综上，本文的研究构建起了基于运行数据的联合循环机组性能分析框架，并揭示了机组运行特性和效率影响因素。后续工作将进一步完善模型细节，深入分析影响机理，并开展面向优化的工况寻优研究，以期为热电联产机组的高效运行提供更全面的理论支持和实践指导。
